\section{Rational numbers}
The procedure parameters and the constants we have used so far
have integer values. In \python, we say that they are of {\em type}\index{Type}
{\idt int} (short for \underline{int}eger). Unlike many
other languages, \python\ imposes no limit on the size of integers: {\idt 2 ** 100}
(2 to the power 100) is $1{,}267{,}650{,}600{,}228{,}229{,}401{,}496{,}703{,}205{,}376$, which leaves
plenty of room for counting the sides of a polygon or the squares of a grid. They are encoded
in base two in the memory of the computer, on as many bits as needed.

In \python\ we can also use decimal numbers with a decimal point such as {\idt 3.14159}, called floating-point numbers (of type \texttt{float}). They
are written according to the English-speaking convention with a point (whereas the French
use a comma). Mathematicians call these numbers {\em rational numbers}\index{Rational}. In \python,
integers are distinguished from rational numbers by the presence of the point. For example
{\idt 1000} is an integer whereas {\idt 1000.0} is a rational number, which can
also be written {\idt 1000.} without any digit after the point. Likewise {\idt .5} is allowed, but
it is more readable to write {\idt 0.5}.

Integers and rational numbers are encoded in different ways in the memory
of the computer. Rational numbers are of type {\idt float}\footnote{\ From
{\em floating-point number}. The computer
keeps in memory a fixed number of significant digits and the position of the decimal point
(which ``floats'' among these digits). These are not real
numbers since numbers of type {\idt float} cannot have infinitely many
non-zero digits after the decimal point!}.

\section{Parameters of type ``float''}
In \python, the type of the parameters of a
procedure is not indicated: the same parameter can receive, from one call to another, an integer or a
rational, and it is only during execution, when the value is used,
that \python\ checks that it is suitable (we say that \python\ is a
{\em dynamically typed} language\index{Dynamic typing}). For example, in the procedure:

\begin{center}
{\mcl def}\ {\idt arc(l):\ gs(l);\ cr()}
\end{center}

\noindent the formal parameter {\idt l} (read the letter $\ell$ and not the digit $1$) is passed to the command {\idt gs} of the \robot\ which
expects a number, integer or rational.

The corresponding actual parameter can therefore be either an integer, as in
{\idt arc(100)}, or a rational number, written with a point, as in {\idt arc(66.666666667)}:
\python\ automatically converts the integer {\idt 100} into the rational {\idt 100.0}
when necessary. If we want to make this conversion ourselves, we use the
function {\idt float}\index{float}: if {\idt n} is an integer,
{\idt float(n)} is the corresponding rational number (for example
{\idt float(3)} is {\idt 3.0}). This is rarely useful.

The command procedures of the \robot\ {\idt gs}, {\idt gsX} and {\idt gsY} (as well as
{\idt gsRT}, {\idt rot} and {\idt fdl} which we will see in the next chapter) accept rational
parameters. This is why we can write {\idt gs(20)} as well as {\idt gs(20.5)}.
The commands {\idt fdn}, {\idt bg}, {\idt pw}, {\idt moveTo} or {\idt drawTo}
on the contrary have integer parameters: {\idt fdn(2.5)} stops the program with the message
`{\tt TypeError:\ 'float' object cannot be interpreted as an integer}'.

\begin{exemple}[Arithmetic spiral]{\rm
The spiral below is made of arcs of circles with successive radii
$100$; $100 / 1.5$ = $66.666\ldots$; $66.666\ldots / 1.5$ = $44.444\ldots$;
$44.444\ldots / 1.5$ = $29.629\ldots$; etc.:

 {\raggedright
 \hbox to 0pt{\hspace*{2cm}\vbox to 0pt{\includevspacePostScript[4cm]{SpiraleArithmetique}{0.8}{0.8}}}
 \input python/ArithmeticSpiral.tex
 }
\noindent $\Box$ }\end{exemple}

\section{Rational expressions}\index{Rational (real) expression}
Addition, subtraction, multiplication and division of rational numbers are
written {\idt +}, {\idt -}, {\idt *} and {\idt /}, that is as for integers. For example
{\idt (((2 * (10.3 + 0.7)) - 1.0) / 3)} is {\idt 7.0}

When an operation involves an integer and a rational, \python\ converts the integer into a
rational and the result is a rational: {\idt 2 * 3.5} is {\idt 7.0}.
When both operands are integers, {\idt +}, {\idt -} and {\idt *} give an integer.
Division, however, requires attention: {\idt /} {\em always} gives a rational,
even between two integers. Thus {\idt 7 / 2} is {\idt 3.5} and {\idt 6 / 2} is {\idt 3.0}.
Integer division, which gives an integer, is written {\idt //}: {\idt 7 // 2} is {\idt 3}.
Thanks to these two notations, the reader of the program always knows which division is performed.
Since {\idt 6 / 2} is a rational, {\idt fdn(6 / 2)} is rejected by the command {\idt fdn}
which expects an integer; one must write {\idt fdn(6 // 2)}.

If $r$ is a number, {\idt sqrt($r$)} is its square root $\sqrt{r}$ (short for
\underline{sq}uare \underline{r}oo\underline{t}),
{\idt $r$ ** 2} is its square $r^{2}$ (which can also be written {\idt $r$ * $r$}),
{\idt abs($r$)} is its \underline{abs}olute value (for an integer as well as for a rational),
{\idt int($r$)} is its integer part obtained by removing the digits after the decimal point
(we say that it is truncated) and {\idt round($r$)} is the nearest integer (from
the verb to \underline{round}). The results of {\idt int} and of {\idt round}
are integers. For example
{\idt int(0.3)} = {\idt int(0.5)} = {\idt int(0.9)} = {\idt 0} whereas
{\idt round(0.3)} = {\idt 0} and {\idt round(0.5)} = {\idt round(0.9)} = {\idt 1}\footnote{\ The
original {\idt round} of \python\ rounds halves to the nearest even integer
({\idt round(0.5)} is {\idt 0} and {\idt round(2.5)} is {\idt 2}). The library of the
\robot\ replaces it by a {\idt round} which rounds halves away from zero, like
Pascal and \ocaml: {\idt round(2.5)} is {\idt 3} and {\idt round(-2.5)} is {\idt -3}.}.
The function {\idt sqrt} is provided by the library of the \robot\ (line
{\mcl from}\ {\idt LibRobot}\ {\mcl import}\ {\idt *}), the others are part of \python.
As for procedures, the argument of a function is always placed between
parentheses: {\idt sqrt(x + 1)}.

In a rational expression, one can use the horizontal length {\idt valGsX()} and the
vertical length {\idt valGsY()} of a square of the grid on which the \robot\ moves. For example
{\idt gsX(valGsX() * 2)} doubles the horizontal size of a square of the grid.

\begin{exemple}[Golden spiral]{\rm
The golden spiral below is made of successive arcs of circles whose radii
have as ratio the {\em golden ratio\/}
$\Phi = \frac{1}{2}(1 + \sqrt{5}) \simeq 1.618033989$. For the advocates of traditional
aesthetics, the golden spiral is said to be a geometric figure very harmonious
to the eye.

{\raggedright
 \vbox to 0pt{\includevspacePostScript[4cm]{SpiraleDor}{0.95}{0.95}}
 \input python/GoldenSpiral.tex
 }
\noindent $\Box$ }\end{exemple}

\section{Rounding errors}
For real numbers\index{Real}, which have infinitely many digits after the decimal point, such as
$\frac{1}{3} = 0.3333333333333333\ldots$ or $\pi = 3.1415926535897932\ldots$
the computer can keep in its memory only a few digits after the decimal point, which
leads, as with a pocket calculator, to rounding errors in computations. For
example in ma\-the\-ma\-tics we have:
\[\left(\left(\left(\left(\left(\sqrt{\sqrt{\sqrt{\sqrt{\sqrt{2}}}}}\right)^{2}\right)^{2}\right)^{2}\right)^{2}\right)^{2} =
\left(\left(\left(\left(\sqrt{\sqrt{\sqrt{\sqrt{2}}}}\right)^{2}\right)^{2}\right)^{2}\right)^{2} =
\] \[\left(\left(\left(\sqrt{\sqrt{\sqrt{2}}}\right)^{2}\right)^{2}\right)^{2} =
\left(\left(\sqrt{\sqrt{2}}\right)^{2}\right)^{2} =
\left(\sqrt{2}\right)^{2} =
2\]
whereas the computer finds that:
\begin{center}
{\idt ((((sqrt(sqrt(sqrt(sqrt(sqrt(2))))) ** 2)\hspace*{6em}}\\
{\idt \hspace*{6em}** 2) ** 2) ** 2) ** 2\ = 1.9999999999999964}
\end{center}

\noindent Likewise, it finds that {\idt 0.1 + 0.2} is
{\idt 0.30000000000000004} (the program {\idt RoundingErrors} of the folder
{\idt book/\allowbreak 19\_\allowbreak rational\_\allowbreak expressions} performs these two computations).

\noindent  Similar rounding problems occur with the grid on which the \robot\ moves. Ideally,
the Cartesian coordinates of the \robot\ should be real numbers.
The computer keeps in memory a rational approximation of them. On the screen they must
be rounded to the nearest integer to correspond to a pixel of the screen.

\needspace{6\baselineskip}\noindent \rule{137mm}{0.1mm}\par\nopagebreak %---------------------------------------------
\begin{exercice}{Write \python\ programs to produce the patterns of the figure
below (only the central pattern, which uses the trig\-o\-no\-met\-ric func\-tions in\-tro\-duced
in the addenda of chapter \ref{ch:rec}, is difficult).
\noindent $\Box$}\end{exercice}
\begin{figure}[htbp]
 \includevspacePostScript[4.6cm]{MotifsDeBeard}{0.95}{0.95}
 \caption{Horemis and Beard patterns; Baravelle spiral}\label{fig:beard}
\end{figure}
\vfill
\begin{corrige}{
We give one program, the others being in the folder
{\idt book/\allowbreak 19\_\allowbreak rational\_\allowbreak expressions}:
\label{prog:CarresImbriquesPeints}
\input python/PaintedNestedSquares.tex
\noindent $\Box$}\end{corrige}

%%%%%%%%%%%%%%%%%%%%%%%%%%%%%%%%%%%%%%%%%%%%%%%%%%%%%%%%%%%%%%%%%%%%%%%%%%%
\noindent \framebox[\textwidth]{
\begin{minipage}{\myparboxwidth}
\setlength{\parindent}{2em}
\small \sf
\vspace{\abovedisplayskip}
\section*{Addenda}
Rational numbers are often called real numbers by abuse of language. They
can be written using {\em scientific notation}\index{Scientific notation}, which
consists in
following the fractional part by a scale factor made of the letter {\idt e} or
{\idt E} (which reads ``ten to the power'') followed by an exponent.

\noindent For example $-3.0${\idt e}$2$ = $-300.0$,
$3.0${\idt e}$0$ = $3.0$ and $3.0${\idt e}$-2$ = $0.03$. A rational number contains no blank. Consequently
$-3.0\hspace{0.3ex}\blanc\hspace{0.3ex}${\idt e}$\hspace{0.3ex}\blanc\hspace{0.3ex}2$ is wrong. In \python,
numbers of type {\idt float} are represented with 15 to 16 significant digits in base 10, the smallest
strictly positive rational (without loss of significant digits) being of the order of
$2.2${\idt e}$-308$ and the largest of the order of $1.8${\idt e}$308$. Beyond that, the result
of a computation is {\idt inf} (infinity): {\idt 1e308 * 10} is {\idt inf}.

\vspace{\abovedisplayskip}
\end{minipage}
}

%%%%%%%%%%%%%%%%%%%%%%%%%%%%%%%%%%%%%%%%%%%%%%%%%%%%%%%%%%%%%%%%%%%%%%%%%%%
%%%%%%%%%%%%%%%%%%%%%%%%%%%%%%%%%%%%%%%%%%%%%%%%%%%%%%%%%%%%%%%%%%%%%%%%%%%
\noindent \framebox[\textwidth]{
\begin{minipage}{\myparboxwidth}
\setlength{\parindent}{2em}
\small \sf
\vspace{\abovedisplayskip}
\smallskip

\index{Rational (real) expression}\index{Integer expression}The value of an arithmetic expression
can be of type {\idt int} or {\idt float}, according to the following rules:
\begin{list}{$\scriptscriptstyle\bullet$}{
}
\item integer constants (without a point) and the operators $+$, $-$, $\ast$, {\idt //},
{\idt \%}, $\ast\ast$ (with a positive exponent) and {\idt abs}\footnote{\ We use different typefaces
for {\em keywords}\index{Reserved identifier} (such as {\mcl def}
or {\mcl for})
and {\em predefined identifiers}\index{Predefined identifier}
(such as {\idt sqrt} or {\idt abs}). The identifiers introduced in a \python\ program
(such as the names of constants, of procedures, etc.) must be
different from the keywords but may be identical to predefined identifiers. This
practice is not recommended because programs which change the usual meaning of predefined
identifiers are unreadable.} applied to integers give an integer result
of type {\idt int}, as do {\idt int} and {\idt round} which
apply to a rational;
\item rational constants (with a point or an exponent {\idt e} or
{\idt E}), division {\idt /} and the function {\idt sqrt} give a rational result
of type {\idt float}, as does {\idt float} which applies to an integer;
\item the operators $+$, $-$, $\ast$, $\ast\ast$ and {\idt abs} applied to rationals,
or to an integer and a rational, give a rational result.
\end{list}
An expression which mixes both, such as {\idt 1 + 2.5}, is therefore allowed: the integer is
converted into a rational and the result {\idt 3.5} is rational.

\smallskip

In \python, the type of parameters is not declared. \python\ automatically converts
an integer into a rational when needed, but never a rational into an integer: it is up to
the programmer to ask for it with {\idt int} or {\idt round}, for example
{\idt fdn(round(x))}. The expected type of a formal parameter can be indicated in the
declaration of the procedure, as in {\mcl def}\ {\idt arc(l:\ float):\ gs(l);\ cr()},
which is not necessary (\python\ does not check it) but can help the reader of the program.

\vspace{\abovedisplayskip}
\end{minipage}
}
%%%%%%%%%%%%%%%%%%%%%%%%%%%%%%%%%%%%%%%%%%%%%%%%%%%%%%%%%%%%%%%%%%%%%%%%%%%
